% Conclusion and Future Work
%
% 1. Summary of review findings (answering RQ1-RQ4 briefly)
% 2. Synthesis of implications
% 3. Research roadmap: prioritized open problems
% 4. Concluding remarks

\subsection{Summary of Findings}
This systematic literature review has synthesized 149 primary
studies (2018--2026) on next-generation software testing, spanning
large-language-model-based test generation, multi-agent architectures,
self-adaptive and self-healing test automation, test oracles, and
autonomous quality engineering. We answered four research questions:

\begin{itemize}
  \item \textbf{RQ1 (Taxonomy):} We identified 18 technique
        categories across the test life cycle, from unit test
        generation to multi-agent systems and self-healing automation.
  \item \textbf{RQ2 (Autonomy):} Most techniques operate at
        human-in-the-loop assistance levels (L1--L3); true agentic
        autonomy (L4--L5) is concentrated in narrow sub-tasks such as
        GUI exploration and flaky-test triage, with significant human
        supervision remaining for test design and oracle validation.
  \item \textbf{RQ3 (Evidence):} Empirical evidence demonstrates
        consistent coverage gains from LLM-based and multi-agent
        generation, but fault-detection effectiveness is less rigorously
        measured, and several widely cited findings do not hold under
        corrected, leakage-free evaluation.
  \item \textbf{RQ4 (Challenges):} Key open problems include the
        test-oracle problem, benchmark contamination,
        cross-language/cross-scale generalization, context management
        for large codebases, and the security and trustworthiness of
        autonomous testing agents.
\end{itemize}

Building on the primary synthesis, this review extended its scope to
address four emergent \emph{quality engineering safety} threads:
security of testing/coding agents, test-oracle independence and the
circularity of AI testing AI-written code, self-healing/flaky-test
stability, and governance/accountability for high levels of autonomy.
This extension was driven by an empirical finding in the corpus itself:
nearly a third of the categorized primary studies address risk and
trust rather than capability. The resulting synthesis formalizes
several of these RQ4 gaps as a connected set of ten open problems and
four candidate, as-yet-unvalidated safety mechanisms.

\subsection{Research Roadmap: Ten Open Problems}
\label{sec:research-roadmap}
Building on the gaps surfaced throughout the synthesis
(Section~\ref{sec:discussion}), we propose the following research
roadmap for next-generation, autonomous software testing. Each item is
stated as an open problem rather than a solution, to orient future
empirical and methodological work (RQ4). Items 1, 4, 7, 8, 9, and 10 —
collectively our \emph{Quality Engineering Safety} thread — are backed
by a dedicated literature synthesis and a candidate safety mechanism each
(detailed in Section~\ref{sec:discussion}; corresponding primary citations
are included in the bibliography):

\begin{enumerate}
  \item \textbf{Oracles independent of the code under test.} Most
        LLM-based oracles are trained or prompted on artifacts derived
        from the same codebase they judge, risking circular validation;
        open problem: oracle construction strategies that do not
        depend on the implementation being tested. Empirical evidence
        of shared generator/judge failure modes is reported by
        \cite{panickssery2024evaluators,wataoka2024selfpreference,
        huang2023selfcorrect,olausson2023selfrepair}; partial mitigations
        (consensus panels, context grounding) are evaluated by
        \cite{xu2025candor,bodicoat2026oracle,molina2024oracleautomation}.
        We propose \emph{Provenance-Partitioned Oracle Verification
        (PPOV)} as a candidate mechanism (formalized in
        Section~\ref{sec:safety-framework}).
  \item \textbf{Probabilistic guarantees on test-suite adequacy.}
        Coverage metrics do not bound fault-detection probability for
        LLM-generated suites; open problem: statistically grounded
        adequacy criteria with confidence guarantees.
  \item \textbf{Contamination-free fault-detection benchmarks.}
        Existing benchmarks risk training-data leakage, inflating
        reported accuracy; open problem: benchmarks that measure true
        fault detection and are verifiably free of contamination.
  \item \textbf{Stability of self-healing loops.} Self-healing locators
        and tests can silently suppress genuine regressions while
        recovering pass rates; open problem: bounding and certifying the
        correctness of automated repairs under distributional UI and
        API drift. Empirical evidence on flaky-classifier limits is
        reported by \cite{alshammari2024flakyclassifiers,parry2023cannier,
        leinen2026detectedundetected}, systemic flakiness by
        \cite{parry2025systemic}, fix stability by \cite{lam2020lifecycle},
        and continuous scoring by \cite{meta2020pfs}. We propose
        \emph{Mutation-Gated Dual-Signal Heal Verification (MG-DSV)} as
        a candidate mechanism (formalized in
        Section~\ref{sec:safety-framework}).
  \item \textbf{Cross-language and cross-ecosystem generalization.}
        Research is heavily concentrated in Java and Python; open
        problem: techniques and benchmarks spanning systems languages,
        functional paradigms, and embedded platforms.
  \item \textbf{Scalable context management for repository-level testing.}
        Single-file/method context limits LLM reasoning; open problem:
        hierarchical retrieval, code-graph representations, and memory
        architectures for whole-repository test generation and triage.
  \item \textbf{Security and integrity of testing agents.} Testing agents
        execute code with elevated privileges, creating attack vectors
        (prompt injection, code poisoning, exfiltration); open problem:
        least-privilege execution environments, verified tool calling,
        and supply-chain defenses tailored to testing pipelines. Attack
        evidence is provided by
        \cite{zhan2024injecagent,debenedetti2024agentdojo,guo2024redcode,
        she2025queryipi,sakib2026trustverify,kozak2025developeraid}, with
        candidate defenses by
        \cite{shi2025progent,wu2024isolategpt,wu2026skillshield,
        xie2026vibeguard,kokkonis2025rosa}. We propose a layered
        \emph{Trust Boundary for Test Agents (TBTA)} architecture
        (formalized in Section~\ref{sec:safety-framework}).
  \item \textbf{Human--agent collaboration.} When to escalate to a
        human, who is accountable for agent-authored tests, and how
        trust is calibrated remain open; open problem: escalation and
        accountability frameworks for mixed human/agent QE teams.
        Evidence on overreliance and trust calibration is reported by
        \cite{passi2025overreliance,ibrahim2025overreliance,ibm2024trew},
        with QA-specific escalation proposals in \cite{kiran2026humanai}.
  \item \textbf{The circularity problem of AI testing AI-written code.}
        As AI increasingly writes both code and its tests, independence
        between implementation and verification erodes; open problem:
        validation strategies that preserve independence when both
        artifacts share a generative origin. See the oracle-
        independence evidence under item 1 in this roadmap
        \cite{panickssery2024evaluators,wataoka2024selfpreference},
        which is the same underlying mechanism.
  \item \textbf{Governance standards for Level-5 (fully autonomous)
        testing.} Existing standards (e.g., ISO/IEC/IEEE~29119) do not
        address fully autonomous test agents; open problem: extending
        or replacing such standards with governance criteria for L5
        autonomy. The current standards baseline is
        \cite{iso2020tr29119,iso2025ts42119}; autonomy-level and
        accountability proposals that could extend it are given by
        \cite{feng2025autonomylevels,treude2026accountable,
        cheng2026threepillar}. We propose an \emph{Autonomy-Tiered
        Quality Sign-off Protocol (ATQSP)} as a candidate extension,
        elaborated in Section~\ref{sec:safety-framework} into the fuller
        Safety-Assured Agentic Quality Engineering (SAAQE) framework,
        with a proposed public-dataset validation protocol in
        Section~\ref{sec:saaqe-validation}.
\end{enumerate}

These ten items map onto, and partially answer, recurring gaps
identified across the included studies (Section~\ref{sec:discussion}):
test-oracle reliability and independence, unmeasured fault-detection
adequacy, benchmark/data contamination, self-healing/flaky-test
stability, reproducibility of agentic results, and security/governance
of autonomous agents.
A cross-thread observation from the Quality Engineering Safety
synthesis: items 1, 4, 7, and 10 share a common structural pattern —
the currently used ``healed / passed / approved'' signal is a single,
gameable proxy that is not independently verified — motivating a
general design principle for next-generation QE safety mechanisms:
\emph{require a second, independently-sourced signal before trusting
the first} (see Section~\ref{sec:qe-safety} and Section~\ref{sec:safety-framework}).
